% =====================================================================
% Literature Review  (~520 words)
% =====================================================================
\section{Background and Literature Review}
\label{sec:literature}

\subsection{Articulated manipulators and degrees of freedom}
A manipulator's degrees of freedom set an upper bound on the classes
of end-effector pose it can reach. Three independent joints are
sufficient to position a point in 3-D space; six are required to also
control orientation~\cite{siciliano2009}. Pick-and-place sorting --- the
task targeted by the autonomous mode of this project --- is well served
by four DOF plus a binary gripper: base rotation, shoulder, elbow, and
a wrist axis to present the gripper correctly to the part. This is
the configuration mandated by the brief and built here.

\subsection{Performance metrics for educational manipulators}
ISO~9283 defines the dominant performance specifications for
manipulating industrial robots, including pose
repeatability and pose accuracy~\cite{iso9283}. Industrial systems
routinely achieve sub-millimetre repeatability, but cost-effective
prototypes built from hobby-grade servos typically perform one to two
orders of magnitude worse without additional sensing. Closing a feedback
loop around each joint is a well-established mitigation that trades a
small computational cost for substantially improved repeatability and
robustness to payload variation~\cite{spong2020,siciliano2009}.

\subsection{Wireless protocols for embedded teleoperation}
The brief permits ``Bluetooth or any wireless communication'', and
the rubric explicitly rewards a justified comparison of options. The
2.4\,GHz ISM band hosts four families of relevance:
Bluetooth Classic (IEEE~802.15.1), Bluetooth Low Energy (BLE), Zigbee
(IEEE~802.15.4) and Wi-Fi (IEEE~802.11), with proprietary GFSK
transceivers (e.g.\ Nordic nRF24-series) representing a fifth
non-standard option~\cite{ieee802151,gomez2012,ieee802154,baker2005}.

Bluetooth Classic with the Serial Port Profile (SPP) provides a
transparent stream-oriented tunnel that emulates an RS-232 link over
the radio --- the mode in which the HC-05 module used in this project
operates. Adaptive frequency hopping across 79~channels of 1\,MHz
gives good resilience against narrowband interference. Reported
round-trip latencies for SPP on commodity modules sit in the
30--100\,ms range~\cite{gomez2012}, which is acceptable for
teleoperation provided the link is permanently established. BLE
offers far lower average current at the cost of connection-interval
quantisation; Zigbee is mesh-capable but high-latency per hop;
Wi-Fi is high-throughput but heavy-stack and slow to associate.
Baker's widely-cited comparison concludes that Bluetooth has the edge
for point-to-point real-time links, while Zigbee's strength is
scalability rather than latency~\cite{baker2005} --- a finding that
directly supports this project's choice of Bluetooth Classic.

\subsection{Microcontroller selection for real-time control}
A Raspberry Pi was considered and rejected. Although a Pi offers
significantly more compute, it runs Linux and therefore cannot
guarantee deterministic timing for hard-real-time PWM and PID loops
without specialist real-time extensions. It also lacks a native ADC,
draws several watts at idle, and undermines the embedded-systems
framing of the brief. By contrast, AVR microcontrollers such as the
ATmega2560 provide deterministic execution, native PWM and ADC
peripherals, multiple hardware UARTs and instant-on
operation~\cite{atmega2560}. The compromise position adopted by this
project, and detailed in Section~\ref{sec:vision}, is to keep the
arm-side controller on the Mega for motor and feedback control, and
introduce a more capable companion processor only for the
computer-vision pipeline.

\subsection{Closed-loop joint control}
\AA str\"om and Murray~\cite{astrom2008} provide the canonical
introduction to feedback control. For servo-driven joints with
potentiometer feedback, a digital PID controller running at a fixed
sample rate (typically 50--200\,Hz) is the standard
choice~\cite{spong2020}. The implementation here samples each
potentiometer at 100\,Hz under timer interrupt and applies a PID
output to the servo's commanded angle. Section~\ref{sec:closedloop}
documents the implementation and tuning.
